% !TeX root = ../../main.tex
\section{Elementi di Teoria della Probabilità}

\epigraph{La teoria della probabilità fornisce strumenti molto generali per il calcolo. Imparare ad usarli con agilità è certamente un’arte, che richiede predisposizione, fantasia, interesse, curiosità, amore per i problemi matematici. Probabilmente pochi sono destinati a diventare artisti, ma tutti possono essere dei buoni artigiani, e questo è quello che conta per il progresso dell’umanità.}{\textit{Sandro Bellini}}

Per iniziare a comprendere il ruolo della teoria della probabilità può essere utile ricordare da quali motivazioni pratiche sia nata, qualche secolo fa. I primi di cui sia documentato l’interesse per questi problemi sono stati giocatori d’azzardo, seguiti dagli assicuratori sulla vita. Fortunatamente la probabilità ha attirato anche l’attenzione di alcuni dei migliori matematici e ha potuto svilupparsi trovando poi numerosissime applicazioni.

Il professionista del gioco d'azzardo è in grado di riconoscere nei risultati di esperimenti casuali (come un lancio di una moneta o l'estrazione di una carta dal mazzo) una certa regolarità che diviene evidente se il numero di prove diventa elevato.

